
\subsubsection{Structural Contracts (Q1): 100\% Detection at Import Time}

Structural contracts successfully detected 2/2 real defects at import time with 0/1 false positives.

\begin{table*}[htbp]
\centering
\resizebox{\dimexpr 2\columnwidth+\columnsep\relax}{!}{%
\begin{tabular}{llllll}
\toprule
Test Scenario & Expected & Actual & Detected & Stage & Time \\
\midrule
Control & 3 channels & 3 channels & N/A & N/A & 0.028s \\
Shape mismatch & 3 channels & 4 channels & ✓ & Import & 0.0003s \\
Dtype mismatch & float32 & float16 & ✓ & Import & 0.0012s \\
\bottomrule
\end{tabular}
}
\end{table*}

\textbf{Detection Rate}: 100\% (2/2 defects)
\textbf{False Positive Rate}: 0\% (0/1 control)
\textbf{Execution Overhead}: <0.03s per validation

\textbf{Key Finding}: Shape mismatches are detectable via layer introspection at import time. The conv layer expecting 4 channels but receiving 3 was flagged in 0.0003s with an actionable error message. Dtype mismatches are detectable via decorator validation in 0.0012s.

\textbf{Interpretation}: These results demonstrate that structural contracts can detect certain classes of API mismatches at import time with minimal overhead. The proof-of-concept validates technical feasibility on two real scenarios.

\subsubsection{Metamorphic Contracts (Q2): 100\% Detection in 3.7ms}

Metamorphic contracts achieved 100\% detection on synthetic violations with zero false positives.

\begin{table*}[htbp]
\centering
\resizebox{\dimexpr 2\columnwidth+\columnsep\relax}{!}{%
\begin{tabular}{lllll}
\toprule
Metric & Baseline & Proposed & Gate Threshold & Status \\
\midrule
**Detection Rate** & 0.0\% & **100.0\%** & $\geq 70$\% & ✅ PASS \\
Softmax Detection & 0/20 (0\%) & 20/20 (100\%) & - & ✅ \\
Dropout Detection & 0/20 (0\%) & 20/20 (100\%) & - & ✅ \\
**False Positive Rate** & 0/10 (0\%) & 0/10 (0\%) & <5\% & ✅ PASS \\
**Execution Time** & 4.26ms & **3.69ms** & $\leq 10s$ & ✅ PASS \\
\bottomrule
\end{tabular}
}
\end{table*}

\textbf{Improvement over Baseline}: +100.0 percentage points

\textbf{Per-Defect-Type Breakdown}:
\begin{itemize}
\item Softmax sum violations: 20/20 detected (100\%)
\item Dropout identity violations: 20/20 detected (100\%)
\end{itemize}

\textbf{Key Finding}: The proposed contracts executed faster than the baseline (3.69ms vs 4.26ms, a -0.57ms overhead) because early defect detection prevented full computation. All 40 injected violations were successfully detected with zero false alarms on 10 control scenarios.

\textbf{Interpretation}: These results demonstrate that metamorphic properties (softmax sum = 1.0, dropout identity in eval mode) are validatable via lightweight probes on synthetic inputs. The 100\% detection rate on programmatically generated violations validates the mechanism. However, the synthetic nature of defects limits generalizability---real-world violations may exhibit more complex failure modes not captured by controlled perturbations.

\subsubsection{Composition Contracts (Q3): 100\% Detection, 0\% False Positives}

Cross-framework validation contracts achieved perfect detection with zero false positives.

\begin{table}[htbp]
\centering
\resizebox{\columnwidth}{!}{%
\begin{tabular}{llll}
\toprule
Metric & Baseline & Proposed & Improvement \\
\midrule
**Detection Rate** & 62.5\% (15/24) & **100.0\% (24/24)** & +37.5pp \\
**False Positive Rate** & 0.0\% (0/6) & **0.0\% (0/6)** & 0pp \\
**Precision** & 100.0\% & **100.0\%** & 0pp \\
**Recall** & 62.5\% & **100.0\%** & +37.5pp \\
**F1 Score** & 0.769 & **1.000** & +0.231 \\
\bottomrule
\end{tabular}
}
\end{table}

\textbf{Validation Time Performance}:
\begin{itemize}
\item Mean: 0.001s (1.0ms)
\item Median: 0.001s (1.0ms)
\item P95: 0.001s (1.1ms)
\item Max: 0.002s (2.0ms)
\end{itemize}

Validation time consistently <1ms, $5000\times$ faster than the 5s threshold.

\textbf{Per-Contract Detection}:
\begin{itemize}
\item Dtype preservation: 5/24 detected (20.8\%)
\item Shape consistency: 9/24 detected (37.5\%)
\item Numerical tolerance: 24/24 detected (100.0\%)
\end{itemize}

\textbf{Confusion Matrix}:
\begin{table}[htbp]
\centering
\resizebox{\columnwidth}{!}{%
\begin{tabular}{lll}
\toprule
 & Detected & Not Detected \\
\midrule
**Defect** & 24 (TP) & 0 (FN) \\
**Valid** & 0 (FP) & 6 (TN) \\
\bottomrule
\end{tabular}
}
\end{table}

\textbf{Key Finding}: The numerical tolerance contract provided comprehensive coverage, catching all defects including those also detectable by dtype and shape contracts. Multiple contracts can detect the same defect (e.g., shape mismatch also causes numerical divergence). The proposed contracts outperformed the end-to-end baseline by 37.5 percentage points in detection rate.

\textbf{Interpretation}: These results demonstrate that cross-framework conversion defects are detectable via multi-level validation (dtype, shape, numerical tolerance). The 100\% detection rate on 30 synthetic test cases validates the mechanism. However, the test setup used $PyTorch\rightarrow PyTorch$ conversion rather than true cross-framework scenarios ($PyTorch\rightarrow TensorFlow, PyTorch\rightarrow ONNX)$, limiting ecological validity.

\subsubsection{Statistical Summary}

\begin{table*}[htbp]
\centering
\resizebox{\dimexpr 2\columnwidth+\columnsep\relax}{!}{%
\begin{tabular}{lllll}
\toprule
Research Question & Metric & Threshold & Result & Status \\
\midrule
**Q1: Structural** & Detection rate & 100\% & 100\% (2/2) & ✅ VALIDATED \\
**Q1: Structural** & False positive rate & 0\% & 0\% (0/1) & ✅ VALIDATED \\
**Q2: Metamorphic** & Detection rate & $\geq 70$\% & 100\% (40/40) & ✅ EXCEEDED \\
**Q2: Metamorphic** & False positive rate & <5\% & 0\% (0/10) & ✅ EXCEEDED \\
**Q2: Metamorphic** & Execution time & $\leq 10s$ & 3.69ms & ✅ EXCEEDED \\
**Q3: Composition** & Detection rate & $\geq 70$\% & 100\% (24/24) & ✅ EXCEEDED \\
**Q3: Composition** & False positive rate & <10\% & 0\% (0/6) & ✅ EXCEEDED \\
**Q3: Composition** & Validation time & <5s & 0.001s & ✅ EXCEEDED \\
\bottomrule
\end{tabular}
}
\end{table*}

All three proof-of-concept validations exceeded their technical thresholds, demonstrating feasibility of the contract mechanisms. However, the small sample sizes (N=2, N=50, N=30) and synthetic defect injection limit generalizability to production settings.
